% Pursuing Stacks (cotes 134-2 to 134-8): archive cote 134-7, site page 74 (PDF page 75), typescript page 516.
% typescript. Draft compared: ps6.tex:3698-3774.
% AI-transcribed by Claude Opus 5.5 (claude-opus-5-5, adaptive thinking, effort high), Message Batches API, batch ps4 (msgbatch_012Aoc1NuskATPsyBeKSYK7J), 2026-11-12; not yet checked by eye.
\scanpage{134-7/74}
\mspage{516}

(14) \qquad $L\pi.(X) \longrightarrow L\hand{\tilde{}}H.(X)$

more precise than the separate homomorphisms
\[ \pi_i(X) \to \hand{\tilde{}}H_i(X) \ , \]
and applying this to an object $\hand{\tilde{}}K(L.)$ and applying (13) again, we should
get a functorial homomorphism\xout{s}

(15) \qquad $L. \to L\hand{\tilde{}}H.(\hand{\tilde{}}K(L\hand{.}))$

in opposite direction from the adjunction morphism

(1\over{\unsure{5}}{6}*) \qquad $L\hand{\tilde{}}H.(\hand{\tilde{}}K(L\hand{.})) \longrightarrow L.$

(p. 509, (5)), the composition of the two being the identity in L. .
In other words, (15) should be a canonical splitting of the natural
adjunction morphism (16). Take for instance $L. =\pi\hand{[}n\hand{]}$, then the first
member of (16) is the Eilenberg-Maclane homology $L\hand{\tilde{}}H.(\pi,n)$, whose first
non vanishing homology group is $\hand{\pi}$ , and (16) is just the canonical
augmentation
\[ L\hand{\tilde{}}H.(\pi,n)\to\pi\hand{[}n\hand{]} \ , \]
and taking the $\mathrm{Ext}^i_{\underline{Z}}$ of both members with $M\hand{[}n\hand{]}$ (M any object in (Ab))
yields the transposed canonical homomorphism

(17) \qquad $\mathrm{Ext}^i_{\hand{Z}}(\pi,M) \longrightarrow H^{\over{\illegible{1}}{n+i}}(\pi,n; M)$ \quad ,

which is an isomorphism for $\over{\illegible{1}}{i=0}$ (whereas for $i\hand{<} \over{\illegible{1}}{0}$ both members are
zero ). For i = 1 , we get an exact sequence (``universal coefficients'')

(18) \qquad $0 \longrightarrow \mathrm{Ext}^1_{\hand{Z}}(\pi,M) \longrightarrow H^{n+1}(\pi,n;M) \longrightarrow \mathrm{Hom}_{\hand{Z}}(\over{\illegible{1}}{H_{n+1}}(\pi,n)\hand{, M)} \hand{\to 0}$ .

A canonical splitting (15) would \over{\unsure{t}}{y}ield\deleted{s} a canonical splitting of this
exact sequence, which definitely looks somehow as ``against nature'' !
Surely, it should'nt be hard to find, for any given inte\xout{r}ger $n \geq 1$,
\xout{\illegible{1}}\inserted{suitable} abelian groups $\hand{\pi}$ and M\hand{,} such that there does not ex\over{\unsure{ast}}{ist} a split-%
ting of (18) stab\xout{\unsure{m}}le under the action of the group
\[ G \ = \ \mathrm{Aut}(\hand{\pi}) \ \mathrm{x} \ \mathrm{Aut}(M) \]
(acting by ``transport de structure'' on the three terms of (18)); presu-%
mably even\hand{,} looking at Aut($\pi$) should be enough. As I am not at all fami-%
liar with Eilenberg-Maclane cohomology, except a little in the case n=1
i.e. for usual group cohomology, I did'nt try to construct an example
for any n, only one for n=1 i.e. in a non simply connected case, which
again is a little less convincing as if it was \inserted{one} for n=2...\xout{Thusxlatix}

Thus let's take n=1 , a vector space of finite dimension over
the prime field $\underline{F}_2$ , and $M = \underline{F}_2$ , in this case standard calculations
\margin{74}
